\section{模型评价与推广}\label{sec:evaluation}\label{sec:limitations}

\subsection{模型优点}

本文的贡献是把一组互相约束的问题做成了可重跑、可反驳的计算实验。
机制内核是一对组合：网络架构由可演化的生成式编码产生，学习被限制为冻结结构上的当代权重扰动、
且不接受其遗传。由此带来四点价值：

\begin{enumerate}
  \item \textbf{把可搜索的对象上移了一层。}被搜索的是生成结构的基因组，
        于是两类结构变异来源同时变得可测量：同源基因组与不同随机种子、单碱基扰动，
        分别对应 H1 与 H2。代价是支撑与符号在训练期冻结，当代无法改变结构。
  \item \textbf{「后天不遗传」由一句引言变为可测量的对照。}结构冻结、只训当代权重，
        于是当代可塑性带来多少与发育结构本身承担多少在机制上被分开，可分别测量。
        若当代扰动显著改善行为而结构完全不变，则当代可塑性足以承担任务；反之则负担主要在发育结构。
  \item \textbf{可用与多样在同一条链上度量。}可用性走闭环生态仿真，多样性走组内图距离与 Mantel 检验；
        这两个判据不依赖跨种子显著性，故在仅有 3 个正式种子的预算下依然可给出结论。
  \item \textbf{输出同源，使结论可被他人复核。}整条生成链确定性推进，
        事件日志、实验指标与专家轨迹三条输出共用同一编号，故每个报告数字都能回溯到运行产物。
\end{enumerate}

\subsection{适用边界}

\begin{enumerate}
  \item \textbf{规模与外推。}实测结构规模为神经元数中位 $37.0$（$27$--$43$，上限 $48$）、
        支撑连接中位 $192.5$ 条（$105$--$289$），比主流任务上 $10^{6}$--$10^{9}$ 参数的网络小若干数量级。
        本平台回答某机制是否存在、幅度多大；该机制在大规模任务上是否仍有效，超出其外推范围。
  \item \textbf{度量与场景。}默认配置下捕获是确定性的，捕食技巧维度未被检验；
        六类细胞类型是功能抽象，结构统计的结论不直接映射到真实细胞类型计数；
        评价场景限于连续二维觅食—避敌。度量上，三个都叫捕获率的量相差约 7 倍、
        分母严重右偏、能量效率为每步非正速率，引用均值时须并列偏态尾与绝对量。
  \item \textbf{统计预算。}正式种子为 3，$n=3$ 时 $3$ 对 $3$ 置换检验双侧最小 $p\approx0.10$，
        跨条件比较以效应量与跨种子波动为准。环境三组中 \texttt{predator\_rich} 区分力明确，
        另两组只有绝对量差别明显（捕获由 $55$ 增至 $100$，或降至 $38$）。
  \item \textbf{训练链的完成度。}演化与学习的联合效果尚未验证；
        学习前后对照、完整消融与多代正式结果正在按计算预算推进。
        外显率也不完全（$0.71$--$0.87$），基因型给出的是架构的概率性偏置。
\end{enumerate}

\subsection{改进方向}

优先级由误差分析给出：先抬高评测回合数（近免费），把模型效应从块间方差里拉出来；
再启用非确定性捕获，使捕食技巧成为可检验维度；随后补齐行为克隆接入与消融、
落盘多代演化正式结果、补出时延与计算量等效率指标。上述任一改进都须重新声明基线口径。

\subsection{推广价值}

本文的可迁移之处有三。其一，闭环生态仿真把「可用与多样」放在同一条链上度量，
该框架可迁移到其他感知—动作耦合的序贯决策任务。其二，把可演化对象上移到基因组、
以发育算子生成结构，为神经演化的生成式路线提供了一条可复现的实验路径。
其三，确定性重放与同源输出使结论可被第三方复核，适合作为教学与可视化平台。
